\documentclass[letterpaper,12pt,oneside]{article}
\usepackage[top=1in, left=1.25in, right=1.25in, bottom=1in]{geometry}

\usepackage[T1]{fontenc}
\usepackage[utf8]{inputenc}
\usepackage[spanish,es-nodecimaldot,es-tabla]{babel}
\usepackage{caption, subcaption} %figuras
\usepackage{graphicx}
\usepackage{array}
\usepackage{tikz}
\usepackage{imakeidx}

\graphicspath{{./figs/}}
\usepackage{setspace}

\begin{document}
%\frontmatter
\begin{titlepage}
		\setlength{\parindent}{0pt} \setlength{\parskip}{0pt}
	
		\begin{center}
			\vfill 
			
			\begin{minipage}{\textwidth}
				
				\newcolumntype{V}{>{\centering\arraybackslash} m{.17\textwidth} }
				\newcolumntype{C}{>{\centering\arraybackslash} m{.56\textwidth} }
				
				\begin{center}
					\includegraphics[width=4.5cm]{UNAM_LOGO2.png}	
				\end{center} 	
				\begin{center}
					\Huge\textbf{Universidad Nacional Autónoma de México}
				\end{center} 
				
			\end{minipage}
				
			\vfill
			
			\begin{minipage}{0.7\textwidth}
				\begin{center}
					\large Ingeniería en ...
				\end{center}
			\end{minipage}
			% 
			\begin{center}
				\vfill
				\large Estructura de Datos y ALgoritmos I
			\end{center}
			
			\begin{center}
				\vspace*{1cm}
				{\huge FORMATO DE REPORTE}\\
				\vspace*{1cm}
				
				\begin{center}
					ALUMNO: \\
					\large 323197294
				\end{center}
				
				\bigskip
				
				Grupo: \\
				\#\\
				Semestre: \\
                2026-II
			\end{center}
			
			\vfill
			
			\begin{center}
				{México,CDMX. Febrero 2026}\\
			\end{center}
			\cleardoublepage
		\end{center}
	\end{titlepage}
 
\tableofcontents
\clearpage
    
%\mainmatter

\section{Introducción}

Este apartado debe abordar los siguientes puntos:

\begin{itemize}
    \item \textbf{Planteamiento del problema.} Se hace una descripción del problema a resolver.
    \item \textbf{Motivación.} Se describe porqué es necesario dar solución al problema.
    \item \textbf{Objetivos.} Lo que se se espera obtener de darle solución al problema.
\end{itemize}

\section{Marco Teórico}

Compilar con en la terminal transforma código fuente en ejecutables binarios [1][2][3]. El comando básico gcc archivo.c crea un ejecutable a.out, mientras que gcc archivo.c -o programa permite nombrar el archivo de salida, realizando el proceso de preprocesamiento, compilación, ensamblaje y enlazado.  
La ejecución de una aplicación CLI generada con gcc depende del sistema operativo, pero generalmente implica compilar el código fuente a un ejecutable y luego invocarlo desde la terminal [4].

CLI o Command Line Interface, es una herramienta que permite interactuar con el sistema operativo, mediante comandos de texto. En el caso de los sistemas Linux, MacOS o bien PowerShell en Windows, los comandos de CLI nos permiten el poder manipular archivos y ejecutar los programas de una manera eficiente [5][6].

En CLI, el comando “cd”, dado su significado “Change Directory”, es utilizado para cambiar de un directorio a otro de una forma más rápida posible [6]. Además, algunos de sus banderas de apoyo son:  
cd -> regresa al directorio anterior  
cd -> va al directorio personal del usuario  
cd /ruta/del/directorio -> accede a una ruta específica  

Por otro lado, ls es un comando de Linux que le permite al usuario hacer una lista de archivos o directorios desde CLI, mostrando su contenido en la línea de comandos [7]. Asimismo, sus banderas de apoyo son:  
-l -> muestra información detallada de los archivos  
-a -> muestra archivos ocultos  
-h -> muestra tamaños legibles para el usuario  

De igual manera, el comando de mover “mv” se utiliza para mover o renombrar directorios y archivos; sin embargo, este no crea una copia del archivo, sino que cambia su ubicación o nombre [8]. Entre algunas de sus banderas de apoyo se encuentran:  
-i -> pide confirmación antes de sobrescribir  
-v -> muestra el proceso realizado  

Por su parte, Remove o rm se usa para borrar archivos o directorios completos de forma permanente [9], usando:  
-r -> elimina directorios y su contenido  
-f -> fuerza la eliminación sin confirmación  
-i -> solicita confirmación antes de borrar  

Asimismo, mkdir es el comando que nos permite crear nuevos directorios o carpetas, ayudando así a organizar estos de manera rápida [10], con ayuda de las banderas de apoyo:  
-p -> crea directorios anidados  
-v -> muestra el directorio creado  

Además, al usar clear, se nos permite limpiar la pantalla de la terminal eliminando los textos previos [11].

Finalmente, compilar con gcc en la terminal transforma código fuente en ejecutables binarios [1][2][3].

\section{Desarrollo}

Es la descripción de la implementación realizada en el lenguaje de programación, así como las pruebas realizadas para obtener los resultados. Es importante resaltar lo siguiente:

\begin{itemize}
    \item No se debe mostrar código, solo describir funciones o la aplicación de los conceptos teóricos.
    \item Las pruebas es describir las entradas ingresadas y las salidas obtenidas.
\end{itemize}

\section{Resultados}

Mediante capturas de pantalla y una breve descripción seguida de la captura se presentan los resultados finales de su aplicación.

\section{Conclusiones}

Se presenta un análisis de los resultados obtenidos, donde se destaca la importancia de la aplicación de los conceptos teóricos para resolver el problema. Es importante resaltar lo siguiente:

\begin{itemize}
    \item No es describir si les gustó la actividad o no.
    \item No es decir que se obtuvo de la práctica.
    \item No es describir lo que fue difícil.
\end{itemize}

\clearpage
\section*{Referencias}

\begin{itemize}
    \item [1] L. Chaparro, “Cómo compilar código C con gcc,” \textit{Medium}, 2021.
    \item [2] L. Chaparro, “Compilation Process with GCC and C Programs,” \textit{Medium}, 2021.
    \item [3] Universidad de la República, “Compilador GCC,” 2019.
    \item [4] YouTube, “Introducción al compilador GCC,” 2020.
    \item [5] Amazon Web Services, “¿Qué es la AWS Command Line Interface?”, 2024.
    \item [6] Arsys, “¿Qué es una CLI o interfaz de línea de comandos?”, 2023.
    \item [7] freeCodeCamp, “El comando Linux ls,” 2023.
    \item [8] GeeksforGeeks, “mv Command in Linux,” 2023.
    \item [9] IONOS, “Comando rm de Linux,” 2023.
    \item [10] IBM, “Creating directories with the mkdir command,” 2022.
    \item [11] GeeksforGeeks, “clear Command in Linux,” 2023.
\end{itemize}

\end{document}